% =====================================================================
% CONCLUSION
% =====================================================================
\section{Conclusion}
\label{sec:conclusion}

We assessed an LLM-backed web application rather than the model inside
it, in a way another engineer can replay and audit, by taking Garak's
definition of a generator seriously and expressing the whole assessment
as plugins registered into Garak's namespace at run time. Garak core is
unmodified, all 28 runs replay through the stock CLI from retained
configurations, and the analysis layer aggregates Garak's evaluation
entries without ever re-judging an outcome.

The substantive results are relative, deliberately so. Data poisoning
was the most exposed category and the one least dependent on the
model's cooperation. Prompt hardening gave a real $5.7\times$
improvement that is nonetheless not a boundary, since a control failing
$8.3\,\%$ of the time against an attacker who pays nothing per attempt
fails with certainty under repetition. Holding the attack corpus fixed
across ten guardrail configurations localised failure to a layer rather
than to a defence's strength, and showed two of ten performing no
better than no guardrail --- the soft system-prompt rule being
instructive, because stating a rule about a secret requires placing the
secret in context. Success concentrated in the techniques that never
name their objective, most sharply where prompts asking for the model's
\emph{rules} extracted the protected secret $5.6\times$ more often than
prompts asking for the secret. Five mislabelled comments created a
targeted classifier backdoor that leaves aggregate accuracy intact and
whose drift is measurable four budget steps before any label changes,
and trust-tagged retrieval eliminated untrusted retrieval entirely.

The methodological result is the one we expect to transfer furthest,
and the ablation is where it is settled rather than asserted. Removing
the ground-truth channel does not degrade the two largest suites by
some margin; it abolishes them. Substituting a stock generic detector
multiplies reported violations by $3.4$ at unchanged recall; consuming
the application's own leak flag instead of cross-tabulating it would
have reported $0.762$ where the independent detector measured $0.000$;
and a four-line harness detail separates a genuine sweep from the same
configuration executed $N$ times with a trend reported over it. Three
further components changed no number at all and are in the table
anyway, because a study reporting only the controls that happened to
fire describes the apparatus it needed rather than the apparatus that
makes the result trustworthy --- the same discipline that keeps the
LLM-as-a-judge a calibrated, opt-in second opinion on which no reported
rate depends. Read across every finding, the controls that held act on
the application pipeline and the controls that failed act on the
model's disposition; the corollary for severity is the single
cross-tenant SQL hit in 21 attempts, negligible as a rate and
categorical as a finding.

Every absolute rate here is a lower bound for one deployment behind one
1B model. Four directions would sharpen the picture without changing
the method: repeating the identical 28-task suite across a model-scale
ladder, for which the apparatus is already parameterised and which
evidence of scale-dependent poisoning susceptibility makes
pressing~\cite[Table~2]{bowen2025scaling}; a multi-turn attack that
first establishes the false premise B6 needs; replacing term-presence
poison scoring with an output-grounding check; and moving the derived
tables that exist only in this paper into the analysis module, since
until then they are the one part of this study reproducible only by
hand.
